% docs/manual/developer/appendices/E-conventions.tex
% 附录 E：CONVENTIONS

\chapter{编码与命名约定}

本附录扩充并整理 \file{src/openbench/CONVENTIONS.md}，给贡献者一份完整的代码规范参考。

\section{命名约定}

\subsection{年份双轨}

\begin{longtable}{l p{4cm} p{6cm}}
  \toprule
  上下文 & 约定 & 示例 \\
  \midrule
  \endhead
  新 schema (\file{config/schema.py}) & \texttt{years: list[int]} & \texttt{[2004, 2010]} \\
  Adapter 输出 (\file{config/adapter.py}) & \texttt{syear} / \texttt{eyear} & \texttt{syear=2004, eyear=2010} \\
  Evaluator 运行时 (\file{core/evaluation.py}) & \texttt{syear} / \texttt{eyear} & 来自 GeneralInfoReader \\
  \bottomrule
\end{longtable}

\textbf{规则}：新代码用 \texttt{years}（list）。Adapter 与 evaluator-facing 运行时结构沿用 \texttt{syear}/\texttt{eyear}，因 evaluator 依赖这两个具体名字。

\subsection{变量名}

\begin{longtable}{l l p{6cm}}
  \toprule
  上下文 & 约定 & 示例 \\
  \midrule
  \endhead
  Standard（目录/config） & 驼峰 + 下划线 & \var{Evapotranspiration}, \var{Latent\_Heat} \\
  NetCDF 内变量名 & 数据集私有 & \texttt{E}, \texttt{Qle}, \texttt{slhf} \\
  Python 标识符 & snake\_case & \texttt{ref\_source}, \texttt{sim\_source} \\
  \bottomrule
\end{longtable}

\subsection{模块命名}

\begin{longtable}{l l p{5cm}}
  \toprule
  上下文 & 约定 & 示例 \\
  \midrule
  \endhead
  新模块 & \texttt{snake\_case.py} & \file{local.py}, \file{cache.py}, \file{schema.py} \\
  迁移兼容模块 & 保留原名 & \file{Mod\_Statistics.py}, \file{Fig\_*.py} \\
  \bottomrule
\end{longtable}

\textbf{规则}：新代码 \texttt{snake\_case}。迁移过来的老文件保留原名以避免破坏内部引用。

\subsection{Config key 命名}

\begin{longtable}{l p{8cm}}
  \toprule
  上下文 & 约定 \\
  \midrule
  \endhead
  新 config (\file{openbench.yaml}) & \texttt{snake\_case} (\texttt{output\_dir}, \texttt{time\_alignment}) \\
  Adapter 输出（给 evaluator/runtime） & 历史 evaluator key (\texttt{compare\_tim\_res}, \texttt{IGBP\_groupby}) \\
  \bottomrule
\end{longtable}

\subsection{分辨率后缀}

Registry 数据集用分辨率后缀：

\begin{itemize}
  \item \texttt{GLEAM\_v4.2a\_LowRes}
  \item \texttt{GLEAM\_v4.2a\_MidRes}
\end{itemize}

写完整名严格匹配；写 base 名（\texttt{GLEAM\_v4.2a}）由 resolver 按 \yamlkey{project.grid_res} 自动选。

\section{类型注解}

\begin{itemize}
  \item 新代码 \textbf{必须}有完整类型注解（函数签名、类属性）
  \item 迁移代码逐步补，不强制
  \item 用 \texttt{Optional[X]} 而非 \texttt{X | None}（兼容 Python 3.10）—— 实际上两者等价，选一致即可
  \item 复杂类型用 \texttt{TypeAlias} 或 dataclass
\end{itemize}

示例：

\begin{minted}{python}
from typing import Optional

def load_config(path: str | Path) -> OpenBenchConfig:
    ...

@dataclass
class ProjectConfig:
    name: str
    output_dir: str
    years: list[int]
    num_cores: Optional[int] = None
\end{minted}

\section{异常处理}

\begin{itemize}
  \item \texttt{ConfigError}：用户配置问题（YAML 错、字段缺、值非法）
  \item \texttt{ValueError} / \texttt{TypeError}：编程错误（开发者写错）
  \item \texttt{OSError} / \texttt{FileNotFoundError}：I/O 错误（用 raise from 链上下文）
  \item 自定义：\file{util/exceptions.py} 中按需添加
\end{itemize}

\subsection{异常消息}

\begin{itemize}
  \item 必须包含可操作的下一步（\textbf{修法}），不只描述问题
  \item 引用 yaml 字段名（用户能定位）
  \item 数字 / 文件路径要拼到消息里
\end{itemize}

示例（好）：

\begin{minted}{python}
raise ConfigError(
    f"project.years start year ({years[0]}) must be <= "
    f"end year ({years[1]})"
)
\end{minted}

示例（坏）：

\begin{minted}{python}
raise ConfigError("invalid years")
\end{minted}

\section{日志级别}

\begin{itemize}
  \item \texttt{DEBUG}：详细流程；仅 \yamlkey{debug_mode: true} 时可见
  \item \texttt{INFO}：阶段切换、关键决定（cache hit、phase start/end）
  \item \texttt{WARNING}：非致命问题（deprecated 字段、低置信度元信息、自动 fallback）
  \item \texttt{ERROR}：失败但程序继续（单 task 失败、单图失败）
  \item \texttt{CRITICAL}：致命，程序无法继续
\end{itemize}

\subsection{logger 命名}

\begin{minted}{python}
import logging

logger = logging.getLogger(__name__)
# __name__ 自动得 'openbench.config.loader' 等
\end{minted}

按模块分 logger，方便用户按子系统调级别。

\section{Docstring}

\begin{itemize}
  \item 公开类与函数 \textbf{必须} 有 docstring
  \item 用 Google 或 numpy 风格（项目内已混用，统一往 Google 风格收）
  \item 至少描述 Args、Returns、Raises
\end{itemize}

示例：

\begin{minted}{python}
def load_config(path: str | Path) -> OpenBenchConfig:
    """Load and validate an openbench.yaml config file.

    Args:
        path: Path to the YAML config file.

    Returns:
        Validated OpenBenchConfig instance.

    Raises:
        ConfigError: If the file is invalid or validation fails.
    """
\end{minted}

\section{测试粒度}

\begin{itemize}
  \item Unit：单函数 / 单类，纯逻辑，跑 < 1ms
  \item Integration：跨模块（loader 加载文件、CLI 调命令），跑 < 1s
  \item Smoke：仅 import 与版本，跑 < 100ms
  \item E2E：完整 evaluation 流程 —— 当前不在 CI（太慢）
\end{itemize}

\section{Commit message}

参考第 11 章。简略：

\begin{itemize}
  \item \texttt{<type>(<scope>): <subject>}
  \item type: feat / fix / docs / refactor / test / chore / build
  \item scope: cli / config / data / runner / gui / dev / user / ops
\end{itemize}

\section{文件组织}

\begin{itemize}
  \item 一个文件 < 600 行；超过 600 行考虑拆分
  \item 一个文件一个主题；不同子系统不要混
  \item 测试文件与源文件路径对应：\file{src/openbench/X/Y.py} → \file{tests/test\_X/test\_Y.py}
\end{itemize}

\textbf{已知例外}：\file{config/adapter.py} 996 行、\file{config/legacy\_processors.py} 909 行、\file{runner/local.py} 906 行；后两个是历史遗留，正在缩。

\section{import 顺序}

ruff 自动整理，但建议遵循：

\begin{enumerate}
  \item 标准库
  \item 第三方
  \item \modname{openbench} 自身
\end{enumerate}

\begin{minted}{python}
"""Module docstring."""
from __future__ import annotations

import json
import logging
from pathlib import Path
from typing import Any

import numpy as np
import yaml

from openbench.config.schema import OpenBenchConfig
from openbench.util.logging_system import setup_logger

logger = logging.getLogger(__name__)
\end{minted}

\section{从 future import}

新文件统一加：

\begin{minted}{python}
from __future__ import annotations
\end{minted}

让所有类型注解延迟求值（PEP 563），允许前向引用（如类内引用自身类型）。

\section{结尾}

如发现这些约定与代码现状有冲突，按"现状是事实"处理：要么改代码，要么改约定。提 PR 让大家一起对齐。
